% Copyright 2026 Benjamin Stanley Frohman (ORCID 0009-0006-7068-3718).
% Licensed under the Apache License, Version 2.0. See LICENSE and NOTICE.
\documentclass[11pt]{article}
\usepackage[margin=1in]{geometry}
\usepackage{amsmath,amssymb,amsthm}
\usepackage{booktabs}
\usepackage{tabularx}
\usepackage{xurl}
\usepackage{enumitem}
\usepackage{xcolor}
\usepackage{hyperref}
\usepackage{listings}
\usepackage{microtype}
\usepackage[T1]{fontenc}
\usepackage{lmodern}

\hypersetup{
  colorlinks=true,
  linkcolor=black,
  citecolor=black,
  urlcolor=blue!50!black,
  pdftitle={Named open hypotheses, gaps T1--T10, and the current Lean implications},
  pdfauthor={Benjamin Stanley Frohman}
}

\lstset{
  basicstyle=\ttfamily\footnotesize,
  breaklines=true,
  breakatwhitespace=true,
  columns=fullflexible,
  keepspaces=true,
  frame=single,
  framesep=3pt,
  xleftmargin=2pt,
  xrightmargin=2pt,
  aboveskip=0.6em,
  belowskip=0.6em,
  showstringspaces=false
}

\theoremstyle{plain}
\newtheorem{theorem}{Theorem}[section]
\newtheorem{lemma}[theorem]{Lemma}
\theoremstyle{definition}
\newtheorem{definition}[theorem]{Definition}
\theoremstyle{remark}
\newtheorem{remark}[theorem]{Remark}

\newcommand{\Tthree}{\mathbb{T}^{3}}
\newcommand{\R}{\mathbb{R}}
\newcommand{\doi}[1]{\href{https://doi.org/#1}{\texttt{#1}}}
\newcommand{\commit}[2]{\texttt{#2}} % repo display handled in text

\title{Identified strain kernel \(A(\theta,\omega)\) and a conditional regularity criterion for 3D Navier--Stokes (\(\kappa=3/2\)): named open hypotheses, gaps T1--T10, and the current Lean implications\\
\large Working notes; not a Clay claim}
\author{Benjamin Stanley Frohman\\
Independent Researcher, Texas, United States\\
ORCID: \href{https://orcid.org/0009-0006-7068-3718}{0009-0006-7068-3718}\\
\texttt{frohmanbenjamin@gmail.com}}
\date{11 November 2026}

\begin{document}
\emergencystretch=3em
\maketitle

\begin{abstract}
This is a working-notes update of \doi{10.5281/zenodo.22835069}. It does \textbf{not} claim the Clay Millennium statement, and it does \textbf{not} derive any estimate that would close a named gap or prove an open hypothesis.

The earlier record recorded (i)~the Biot--Savart identification of the strain kernel \(A(\theta,\omega)=(\theta\times\omega)\otimes\theta+\theta\otimes(\theta\times\omega)\), (ii)~the prefactor lock \(\kappa=C_{\mathrm{CZ}}(3)=(3/8\pi)\cdot 4\pi=3/2\neq 1\), (iii)~same-point cancellation and the \(|\sin\varphi|\) pairing, and (iv)~a conditional regularity criterion: if a data-only or energy-controlled bound on \(\|S\|_{\infty}\) exists in a space that controls the Beale--Kato--Majda integral, then the classical solution on \(\Tthree\) exists globally. Named gaps T1--T10 (Lemmas~4.1--4.6 in that record's abstract) record why that bound is not derived there. An \(L^{\infty}\) bound on \(S\) on the whole torus is emptiness of the singular set, not \(\mathbb{P}^{1}\)-nullity (Caffarelli--Kohn--Nirenberg). The two-point strain pairing that would discharge Constantin--Fefferman is not free.

This update restates those gaps faithfully from the published PDF, records how the abstract's phrase ``T1--T10'' sits relative to the lemmas actually printed in the body, and inserts the current exact Lean statements of the named open hypotheses and of the conditional implications
\[
\text{NamedHypothesis}\implies\text{GlobalRegularity},
\]
pinned to public commits. The rate and bound remain open. By the Beale--Kato--Majda criterion, the remaining periodic closing bound is as hard as Fefferman's Statement~(B) itself.
\end{abstract}

\noindent
\fbox{\parbox{\dimexpr\linewidth-2\fboxsep-2\fboxrule}{\textbf{This is not a Clay Millennium claim.} The open cubic rate \(M'\le C M^{2}-\kappa M^{3}\) and the uniform vorticity closing bound are unproved for general smooth, divergence-free, finite-energy data. The architecture of the formalizations is an implication from a named \texttt{Prop} to global regularity. Composition would give an unconditional theorem only after a proof of that \texttt{Prop} exists. No such proof is offered here.}}

\tableofcontents
\newpage

\section{Purpose and non-claims}

This note has one job: to put in one place, without promotional language,
\begin{enumerate}[leftmargin=1.4em]
\item the precise statements of the named open Lean hypotheses now used in the author's public repositories,
\item a faithful restatement of the obstruction lemmas of \doi{10.5281/zenodo.22835069} (the gaps document), including the abstract's wording on T1--T10 and the sentence ``the pairing is not free,''
\item the current conditional theorems as implications \(\text{NamedHypothesis}\to\text{GlobalRegularity}\), with exact names and file locations,
\item the honest scope: the rate and bound are open; the remaining periodic bound is equivalent in difficulty to Clay~(B); this is not a Clay claim; and the \(\Tthree=\R^{3}\) caveats and known refutations are recorded rather than smoothed over.
\end{enumerate}

Nothing in this note invents a derivation that would close T1--T10, prove \(M'\le C M^{2}-\kappa M^{3}\) for general data, or convert any conditional theorem into Statement~(B). Line breaks inside Lean listings are typographic only; the repository text at the cited commit and \texttt{file:line} is authoritative.

\section{Choice of record and overlap}

This document is prepared as a \textbf{new version of the existing concept} \doi{10.5281/zenodo.22835068}, whose only published version is \doi{10.5281/zenodo.22835069} (18~September 2026; one file, the official proposed theorem PDF). That record is the author's gaps document: its description is the source of the phrase ``Named gaps T1--T10 (Lemmas~4.1--4.6)'' and of the sentence ``The two-point strain pairing that would discharge Constantin--Fefferman is not free.'' A later GitHub--Zenodo snapshot of the same repository, \doi{10.5281/zenodo.23867482} (concept \doi{10.5281/zenodo.23867481}), is a source zip with a changelog and does not replace the gaps PDF.

A new concept would have duplicated that ledger. The other current records already have their own jobs and are cited, not re-released:
\begin{itemize}[leftmargin=1.4em]
\item \doi{10.5281/zenodo.22240168} --- geometric reduction for the vorticity direction; remaining estimates left open; not a Clay solution.
\item \doi{10.5281/zenodo.24144846} --- Fefferman formalization v1.5.0 (one-sided depletion hypothesis; the old two-sided form is kernel-refuted).
\item \doi{10.5281/zenodo.24144852} (concept \doi{10.5281/zenodo.23749844}) --- conditional reduction note, with the word ``unconditional'' removed and a vacuity erratum added.
\item \doi{10.5281/zenodo.24144853} (concept \doi{10.5281/zenodo.23627609}) --- \texttt{NS\_\hspace{0pt}Millennium\_\hspace{0pt}Proof} v2.5.0, conditional title and explicit ``not a Clay Millennium claim'' scope.
\item \doi{10.5281/zenodo.24144856} (concept \doi{10.5281/zenodo.23652842}) --- tether regularity v1.3.0, with the P5 overclaim withdrawn and the statement that its \texttt{T3} is really \(\R^{3}\).
\item Concept \doi{10.5281/zenodo.24081999} (latest published version \doi{10.5281/zenodo.24122596}, 10~November 2026) --- sdiff formalization; the BKM pillar merge of 11~November 2026 at commit \texttt{4aba751} is on GitHub and is recorded here, but is not itself a new Zenodo version of that concept.
\item \doi{10.5281/zenodo.24066255} --- two-point strain--direction pairing note.
\end{itemize}
The unpublished draft 23940666 is not modified or reused. This note does not cite \doi{10.5281/zenodo.23717084}.

\section{The official conditional criterion of the gaps record}

The published PDF of \doi{10.5281/zenodo.22835069} states the official proposed theorem as follows (Theorem~2.1, notation unchanged in substance). Let \(s>5/2\) and let \(u_{0}\in H^{s}(\Tthree;\R^{3})\) satisfy \(\operatorname{div}u_{0}=0\). Let \(u\) be the unique Kato solution of incompressible Navier--Stokes
\[
\partial_{t}u+(u\cdot\nabla)u+\nabla p=\nu\Delta u,\qquad
\operatorname{div}u=0,\qquad u(\cdot,0)=u_{0},
\]
on its maximal classical interval \([0,T_{\ast})\), with vorticity \(\omega=\operatorname{curl}u\). Assume Lemmas~4.1, 4.2, 4.3, 4.4, 4.5, and 4.6 of that note, with constants depending only on \(\|u_{0}\|_{H^{s}}\) and \(\nu\). Then
\[
\int_{0}^{T_{\ast}}\|\omega(\cdot,t)\|_{L^{\infty}}\,dt<\infty,
\]
hence \(T_{\ast}=\infty\) by the Beale--Kato--Majda continuation lemma cited there, and \(u\) remains smooth for all \(t\ge 0\).

\begin{remark}[Status, quoted in substance from Remark~2.2]
Theorem~2.1 is not the Clay theorem until every cited lemma is a theorem. The lemmas marked ``classical, filled'' in that note are theorems. The lemmas marked ``obstruction'' are not. Filling those hypotheses with first-principles, non-circular, solution-independent estimates would be the Clay proof. Declaring them filled because the names coincide is circular.
\end{remark}

\section{Gaps T1--T10, restated from \doi{10.5281/zenodo.22835069}}

\subsection{How the label ``T1--T10'' is used in the source}

The \textbf{body} of the published PDF never prints the tokens T1,~\dots, T10. The obstruction section is written as Lemmas~4.1--4.8. The abstract and the deposit description say: ``Named gaps T1--T10 (Lemmas~4.1--4.6) record why that bound is not derived here.'' The same abstract adds two scope statements that are not numbered as lemmas: an \(L^{\infty}\) bound on \(S\) on the whole torus is emptiness of the singular set, not \(\mathbb{P}^{1}\)-nullity (Caffarelli--Kohn--Nirenberg); and ``The two-point strain pairing that would discharge Constantin--Fefferman is not free.''

This note does not invent a numbering that the PDF does not print. It restates every named obstruction in that record, and records the following reading of the abstract's phrase T1--T10, which is a ledger reading rather than a typographic fact of the PDF:
\begin{center}
\small
\begin{tabularx}{\linewidth}{@{}>{\raggedright\arraybackslash}p{0.7in}>{\raggedright\arraybackslash}p{1.7in}X@{}}
\toprule
Label & Source in 22835069 & What it blocks \\
\midrule
T1 & Lemma~4.1, strain endpoint & a solution-independent space \(X\) and constant \(C_{\ast}\) with \(\|S\|_{L^{\infty}}\le C_{\ast}\|\omega\|_{X}\) \\
T2 & Lemma~4.2, stretching identification & \(\int\langle\omega_{\varepsilon},(\omega_{\varepsilon}\cdot\nabla)u_{\varepsilon}\rangle\phi_{\varepsilon}\le C_{\ast}M_{\varepsilon}I_{4}(\phi_{\varepsilon})+C_{\varepsilon}\) with commutators compatible with a data-only majorant \\
T3 & Lemma~4.3, transported weight & a weight \(\phi_{\varepsilon}\) whose \(L^{\infty}\) ODE coefficients depend only on \(\|u_{0}\|_{H^{s}}\) and \(\nu\) (if \(\phi\equiv 1\), the tether is ornamental and the problem moves onto Lemma~4.2) \\
T4 & Lemma~4.4, absorption constant & \(C_{\mathrm{abs}}\) depending only on \(\kappa=3/2\), \(C_{\mathrm{Sob}}\), \(\lambda(\Tthree)\), and \(\nu\) \\
T5 & Lemma~4.5, non-circular majorant & a comparison function \(Y\) whose ODE coefficients are independent of \(\{M(t):t\ge 0\}\), with \(M\le Y\) \\
T6 & Lemma~4.6, mollifier limit & commutator constants that remain controlled as \(T\uparrow T_{\ast}\), depending only on \(u_{0}\) and \(\nu\) \\
T7 & Lemma~4.7, Jacobi/cocycle of \(B\) & verification of a tethered bracket that is simultaneously degenerate on \(H\), Jacobi, and a source of a new negative term in \(\frac{d}{dt}M\) (in tension; vacated if \(B\) is dropped) \\
T8 & Lemma~4.8, degeneracy versus feedback & the dichotomy: a Poisson correction with \(B(\cdot,H)=0\) cannot supply a new term in the vorticity equation; a nonzero correction changes the PDE \\
T9 & abstract / description & an \(L^{\infty}\) bound on \(S\) on the whole torus is emptiness of the singular set, not \(\mathbb{P}^{1}\)-nullity (CKN) \\
T10 & abstract / description & the two-point strain pairing that would discharge Constantin--Fefferman is not free \\
\bottomrule
\end{tabularx}
\end{center}

T1--T6 are exactly the six hypotheses of Theorem~2.1. Section~6 of the source says that those five analytic pieces (4.1, 4.2, 4.4, 4.5, 4.6), with 4.3 the weight, ``are the same missing a priori bound under different names.'' Proving them would prove Theorem~2.1 unconditionally, which is the Clay statement. No rearrangement of the names discharges the estimate.

\subsection{Lemma~4.1 (strain endpoint) --- T1}

There exist a Banach space \(X(\Tthree)\) and a constant \(C_{\ast}\) such that \(\|S\|_{L^{\infty}}\le C_{\ast}\|\omega\|_{X}\) for every smooth divergence-free \(\omega\), with \(S\) the deformation tensor of \(u=K_{\mathrm{BS}}\ast\omega\).

The source writes the kernel representation with angular factor \(A(\theta,\omega)=(\theta\times\omega)\otimes\theta+\theta\otimes(\theta\times\omega)\), and records the scale tag \(C_{\mathrm{CZ}}(3)=3/2\). It then stops: the Calder\'on--Zygmund operator \(\omega\mapsto S\) is bounded on \(L^{p}(\Tthree)\) for \(1<p<\infty\) and is \textbf{not} bounded on \(L^{\infty}\). The estimate \(\|S\|_{L^{\infty}}\le\tfrac{3}{2}\|\omega\|_{L^{\infty}}\) is therefore false as a general statement. A Kozono--Taniuchi or \(\dot B^{0}_{\infty,1}\) substitute changes the Young bookkeeping. Chat fragments that set \(C_{\ast}=3/2\) on \(L^{\infty}\) are not used. Until the endpoint is fixed, Lemma~4.2 cannot feed the Young absorption.

\subsection{Lemma~4.2 (stretching identification) --- T2}

On a classical solution, after mollification,
\[
\int_{\Tthree}\langle\omega_{\varepsilon},(\omega_{\varepsilon}\cdot\nabla)u_{\varepsilon}\rangle\phi_{\varepsilon}\,d\lambda
\le C_{\ast} M_{\varepsilon} I_{4}(\phi_{\varepsilon})+C_{\varepsilon},
\]
with \(C_{\varepsilon}\to 0\) in the sense used to pass to the limit in \(S_{\varepsilon}\), and with \(C_{\ast}\) the constant of Lemma~4.1. The algebraic identity \(\langle\omega,(\omega\cdot\nabla)u\rangle=\langle\omega,S\omega\rangle\) is filled. The missing piece is the same \(L^{\infty}\) bound, plus mollifier commutators whose constants may depend on \(\|u\|_{C^{2}([0,T]\times\Tthree)}\) and therefore cannot be fed into a majorant whose coefficients may depend only on \(u_{0}\) and \(\nu\).

\subsection{Lemma~4.3 (transported weight) --- T3}

There exists a family of nonnegative scalars \(\phi_{\varepsilon}\) such that \(\|\phi_{\varepsilon}(t)\|_{L^{\infty}}\) is dominated by a function of \(t\) whose ODE coefficients depend only on \(\|u_{0}\|_{H^{s}}\) and \(\nu\). Linear transport with \(\operatorname{div}u=0\) conserves \(\|\phi\|_{L^{\infty}}\) and is filled. If the weight grows in response to stretching, the ODE inherits the unknown \(C_{\ast}\). Taking \(\phi\equiv 1\) makes the tether ornamental and moves the problem onto Lemma~4.2. Either choice must be declared; they cannot be mixed.

\subsection{Lemma~4.4 (absorption constant) --- T4}

The constant \(C_{\mathrm{abs}}\) in the schematic inequality for \(\frac{d}{dt}S_{\varepsilon}\) depends only on \(\kappa=3/2\), \(C_{\mathrm{Sob}}\), \(\lambda(\Tthree)\), and \(\nu\). Given a stretching contribution of size \(C_{\ast}MI_{4}(\phi)\), Young produces an explicit \(C_{\mathrm{abs}}(C_{\ast},C_{\mathrm{Sob}},\lambda(\Tthree))\). Until Lemma~4.1 supplies a solution-independent \(C_{\ast}\), \(C_{\mathrm{abs}}\) is not a function of \(u_{0}\) and \(\nu\) alone.

\subsection{Lemma~4.5 (non-circular majorant) --- T5}

There exists a scalar \(Y\in C^{1}([0,\infty);[0,\infty))\) whose ODE \(Y'=F(Y)\) has coefficients independent of \(\{M(t):t\ge 0\}\), such that \(M(t)\le Y(t)\) on the existence interval. Scalar ODE comparison, including a clipped cubic \(F(y)=C_{\mathrm{abs}}y^{2}-\kappa'' y^{3}\) with a global barrier, is filled. The missing piece is the differential inequality \(M'\le F(M)\) with an \(F\) built only from \(u_{0}\) and \(\nu\). That estimate is the analytic core of the Clay problem. Inserting \(\sup M\) into \(C_{\mathrm{abs}}\) is circular; modifying the PDE is inadmissible.

\subsection{Lemma~4.6 (mollifier limit) --- T6}

\(S_{\varepsilon}(t)\to S(t)\) and the commutators \(C_{\varepsilon},R_{\varepsilon}\) vanish as \(\varepsilon\to 0\), with constants compatible with Lemma~4.5. Passage on a compact classical interval \([0,T]\subset[0,T_{\ast})\) is filled. Using those bounds to continue past a putative singularity requires constants that remain controlled as \(T\uparrow T_{\ast}\). If they explode with \(M(t)\), they cannot close Lemma~4.5.

\subsection{Lemma~4.7 (Jacobi / cocycle) --- T7}

If a bilinear correction \(B\) is part of the argument, then \(d^{2}B=0\) on the reduced coadjoint orbit. The algebraic expansion is a finite computation; the missing piece is verification that the particular kernel cancels in every slot while remaining degenerate on \(H\) (Lemma~4.8) and contributing a nonzero negative term to \(\frac{d}{dt}M\) (Lemma~4.2). Those three demands are in tension. If \(B\) is dropped and \(\phi\) is an ordinary weight, this lemma is vacated and Theorem~2.1 is a Lyapunov criterion on unmodified Navier--Stokes.

\subsection{Lemma~4.8 (degeneracy versus feedback) --- T8}

A correction used as dynamics satisfies \(B(F,H)=0\). A weight used as a Lyapunov device is not required to be a Poisson correction. The two readings are not mixed. If \(B(\cdot,H)=0\), the Hamiltonian evolution is exactly Euler and no extra term appears in the vorticity equation. If \(B(\cdot,H)\neq 0\), the evolution is not Navier--Stokes and Theorem~2.1 does not apply. This is filled as logic and left open as geometry.

\subsection{The two abstract scope statements --- T9 and T10}

An \(L^{\infty}\) bound on \(S\) on the whole torus is emptiness of the singular set, not \(\mathbb{P}^{1}\)-nullity in the sense of Caffarelli--Kohn--Nirenberg. The two-point strain pairing that would discharge Constantin--Fefferman is not free. Neither sentence is a proof of a bound, and neither is weakened here.

\section{Named Lean hypotheses, copied from the public commits}

All three repositories below use toolchain \texttt{leanprover/\hspace{0pt}lean4:v4.\hspace{0pt}30.\hspace{0pt}0-rc1} and Mathlib commit \texttt{0c154d67103f74be3a0f2c509f72ccbf5be9f2a7}. The statements are \texttt{Prop}s passed as hypotheses, not axioms and not theorems about general Navier--Stokes solutions, except where a refutation or a zero-datum special case is explicitly named.

\subsection{Periodic closing bound (sdiff, commit \texttt{4aba751})}

This is the remaining hypothesis of the sdiff main theorem after the BKM pillar of pull request~\#5.

\textit{Source:} \texttt{BenFrohman/\hspace{0pt}frohmanian-tether-sdiff}, commit \texttt{4aba751}, file \texttt{FrohmanianTether/\hspace{0pt}ClayStatements.\hspace{0pt}lean}, lines 575--581.
\begin{lstlisting}
@[expose] public def PeriodicTetherClosingBound (nu : Real) : Prop :=
  forall u0 : VelocityField,
    (forall x, div u0 x = 0) -> ContDiff Real (top : NatInfinity) u0 ->
      IsZ3Periodic u0 ->
      exists Y : Real, forall (T : Real) (u : TimeDependentVelocity)
        (p : TimeDependentPressure),
        IsPeriodicSmoothSolutionOn nu T u0 u p ->
          forall t in Set.Ico (0 : Real) T,
            vorticity_sup_norm (vorticity (u t)) <= Y
\end{lstlisting}
(The listing spells Lean's unicode as ASCII only in this paragraph's display of operators; the file uses \(\nu\), \(\forall\), \(\to\), \(\top\), \(\exists\), \(\in\), \(\le\).) In words: for every smooth, divergence-free, \(\mathbb{Z}^{3}\)-periodic initial datum, there is a finite ceiling \(Y\), independent of the existence time \(T\), that bounds \(\|\omega(t)\|_{\infty}\) along every classical periodic solution on \([0,T)\).

The pointwise bound for a single datum is \texttt{PeriodicTetherBound} (\texttt{FrohmanianTether/\hspace{0pt}TetherImpliesBKM.\hspace{0pt}lean}, line~45). The \(\R^{3}\)/physical-class companions, which are \textbf{not} the Clay~(B) hypotheses, are recorded next.

\subsection{Physical-class rate and bound (sdiff, same commit)}

\textit{Source:} \texttt{FrohmanianTether/\hspace{0pt}SymplecticTether.\hspace{0pt}lean}, lines 395--401 and 415--425.
\begin{lstlisting}
@[expose] public abbrev FrohmanianTetherBound (u0 : VelocityField) (nu : Real) : Prop :=
  forall u : Real -> VelocityField, forall p : Real -> PressureField,
    NS_PDE u p nu -> u 0 = u0 -> IsPhysicalNSSolution u ->
      exists Y : Real,
        (forall tau >= (0 : Real),
          vorticity_sup_norm (vorticity (u tau)) <= Y) /\
        Continuous (fun tau : Real => vorticity_sup_norm (vorticity (u tau))) /\
        (forall tau >= (0 : Real),
          BddAbove (Set.range fun x => ||vorticity (u tau) x||))

@[expose] public abbrev FrohmanianTetherRate (u0 : VelocityField) (nu : Real) : Prop :=
  forall u : Real -> VelocityField, forall p : Real -> PressureField,
    0 < nu -> NS_PDE u p nu -> u 0 = u0 -> IsPhysicalNSSolution u ->
      exists C k : Real, 0 < C /\ 0 < k /\
        (forall s >= (0 : Real),
          DifferentiableAt Real
            (fun tau => vorticity_sup_norm (vorticity (u tau))) s) /\
        (forall s >= (0 : Real),
          deriv (fun tau => vorticity_sup_norm (vorticity (u tau))) s <=
            C * vorticity_sup_norm (vorticity (u s)) ^ 2 -
              k * vorticity_sup_norm (vorticity (u s)) ^ 3) /\
        Continuous (fun tau : Real => vorticity_sup_norm (vorticity (u tau))) /\
        (forall tau >= (0 : Real),
          BddAbove (Set.range fun x => ||vorticity (u tau) x||))
\end{lstlisting}
The closing form \texttt{FrohmanianTetherClosingBound} (\texttt{ClayStatements.\hspace{0pt}lean}, line~48) asks the bound of every smooth, divergence-free, finite-energy datum in this physical class. Related compact-interval forms are \texttt{LocalTetherRate}, \texttt{LocalTetherBound}, \texttt{LocalTetherClosingBound}, and the earlier two-sided variants \texttt{LocalTetherRateV014} and \texttt{LocalTetherBoundV014}.

\textbf{These physical-class statements are not theorems for general data.} See~\S\ref{sec:refutations}.

\subsection{Repaired rate on Fefferman's smooth class (sdiff, pull request~\#4)}
\label{sec:rate-smooth}

This material entered the public tree in the merge of pull request~\#4 (commit \texttt{0698d08}, 11~November 2026) and is present unchanged on main at commit \texttt{4aba751}. File \texttt{FrohmanianTether/\hspace{0pt}FeffermanClayA.\hspace{0pt}lean}. It is the repaired rate intended for Fefferman's smooth class, after the unrestricted physical-class rate was refuted (see~\S\ref{sec:refutations}).

\textit{Source:} \texttt{BenFrohman/\hspace{0pt}frohmanian-tether-sdiff}, commit \texttt{4aba751}, lines 81--88.
\begin{lstlisting}
@[expose] public def FrohmanianTetherRateSmooth
    (u0 : VelocityField) (nu : Real) : Prop :=
  exists C kap : Real, 0 < C /\ 0 < kap /\
    forall (T : Real) (u : TimeDependentVelocity)
      (p : TimeDependentPressure),
      IsFeffermanSolutionOn nu T u0 u p ->
        ContinuousOn (Momega u) (Set.Ico (0 : Real) T) /\
        (forall s in Set.Ico (0 : Real) T, exists d : Real,
          HasDerivWithinAt (Momega u) d (Set.Ici s) s /\
            d <= C * Momega u s ^ 2 - kap * Momega u s ^ 3) /\
        (forall t in Set.Ico (0 : Real) T,
          BddAbove (Set.range fun x => ||vorticity (u t) x||))
\end{lstlisting}
(The listing spells Lean's unicode as ASCII; the file uses \(\nu\), \(u_{0}\), \(\forall\), \(\exists\), \(\to\), \(\wedge\), \(\in\), \(\le\), \(\|\cdot\|\), and the identifier \texttt{M\(\omega\)}.) In words: along every local solution in Fefferman's smooth class on \(\R^{3}\times[0,T)\) from \(u_{0}\), \(M=\|\omega\|_{\infty}\) is finite, continuous on \([0,T)\), and has a right derivative satisfying \(M'\le C M^{2}-\kappa M^{3}\) with \(C,\kappa>0\) depending only on \(u_{0},\nu\). The solution class \texttt{IsFeffermanSolutionOn} (same file, lines 63--72) requires joint smoothness, the Navier--Stokes equations with \(f\equiv 0\), the initial condition, and the energy inequality. The junk-derivative counterexample of \S\ref{sec:refutations}, whose datum is nowhere differentiable, is excluded.

This \texttt{Prop} is \textbf{open} for general Fefferman data. No derivation of it from the unmodified vorticity equation appears in the formalization or in this note. It is shown only at the zero field: \texttt{frohmanianTetherRateSmooth\_\hspace{0pt}zero} (line~176).

The companion continuation hypothesis on \(\R^{3}\) is \textbf{not} the periodic BKM theorem of pull request~\#5.

\newpage
\textit{Source:} same file, lines 94--102.
\begin{lstlisting}
@[expose] public def R3BKMContinuationAt (nu : Real)
    (u0 : VelocityField) : Prop :=
  (exists Y : Real, forall (T : Real) (u : TimeDependentVelocity)
      (p : TimeDependentPressure),
      IsFeffermanSolutionOn nu T u0 u p ->
        forall t in Set.Ico (0 : Real) T,
          BddAbove (Set.range fun x => ||vorticity (u t) x||) /\
            Momega u t <= Y) ->
    exists (u : TimeDependentVelocity) (p : TimeDependentPressure),
      FeffermanGlobalSolution nu u0 u p

@[expose] public def R3BKMContinuation (nu : Real) : Prop :=
  forall u0 : VelocityField,
    ContDiff Real (top : NatInfinity) u0 ->
      (forall x, div u0 x = 0) -> FeffermanDecay u0 ->
        R3BKMContinuationAt nu u0
\end{lstlisting}
The file's docstring records that this is the classical Kato / Beale--Kato--Majda continuation on \(\R^{3}\), ``known, NOT formalized.'' It is a named hypothesis, not a theorem, for general data. The only kernel instance is \texttt{r3BKMContinuationAt\_\hspace{0pt}zero} (line~146), at \(u_{0}=0\). Do not confuse \texttt{R3BKMContinuation} with \texttt{PeriodicBKMContinuation}, which is proved for every \(\nu>0\) as \texttt{periodicBKMContinuation\_\hspace{0pt}of\_\hspace{0pt}classical} (\texttt{FrohmanianTether/\hspace{0pt}Pillars/\hspace{0pt}BKMContinuation.\hspace{0pt}lean}, line~144).

A global-solution form \texttt{FrohmanianTetherRateSmooth\allowbreak Global} (lines 195--205) is recorded only to be set aside: \texttt{frohmanianTetherRateSmoothGlobal\_\hspace{0pt}of\_\hspace{0pt}no\_\hspace{0pt}global\_\hspace{0pt}solution} (lines 207--210) shows that it holds automatically for any datum that has no global smooth physical solution, and therefore cannot feed Clay~(A).

\subsection{One-sided direction depletion (fefferman, main after v1.5.0)}

\textit{Source:} \texttt{BenFrohman/\hspace{0pt}fefferman-clay-navier-stokes}, commit \texttt{3512caa} (main after the v1.5.0 erratum). The one-sided correction and the kernel refutation of the old form were proved in commit \texttt{afaab80}. File \texttt{FeffermanClayNS/\hspace{0pt}HighSetDepletionAndBKM.\hspace{0pt}lean}, lines 234--246 (current form) and 251--263 (old form).
\begin{lstlisting}
@[expose] public def PeriodicDirectionDepletionBound (nu : Real) : Prop :=
  forall u0 : VelocityField,
    IsPeriodicSmoothInitialData u0 ->
    exists C_CZ k : Real, 0 < C_CZ /\ 0 < k /\
      forall (T : Real) (u : TimeDependentVelocity)
        (p : TimeDependentPressure),
        IsPeriodicSmoothSolutionOn nu u0 (fun _ _ => 0) T u p ->
        ContinuousOn (fun s => vorticity_sup_norm (vorticity (u s)))
          (Set.Icc (0 : Real) T) /\
        (forall s in Set.Ico (0 : Real) T, exists d : Real,
          HasDerivWithinAt
            (fun tau => vorticity_sup_norm (vorticity (u tau)))
            d (Set.Ici s) s /\
            d <= C_CZ * (vorticity_sup_norm (vorticity (u s))) ^ 2 -
              k * (vorticity_sup_norm (vorticity (u s))) ^ 3)
\end{lstlisting}
This is the open cubic rate \(M'\le C_{\mathrm{CZ}}M^{2}-\kappa''M^{3}\) along every local periodic smooth solution, with \(M\) continuous on \([0,T]\) and \(M'\) a \textbf{right} derivative at each \(s\in[0,T)\). The alias \texttt{DirectionCurvatureBound} is this same definition (\texttt{FeffermanClayNS/\hspace{0pt}MainTheorem.\hspace{0pt}lean}, line~32). The hypothesis is \textbf{open}.

The v1.0--v1.4 form, retained as \texttt{PeriodicDirectionDepletionBoundOriginal}, asked for a two-sided \texttt{DifferentiableAt} at every \(t\in[0,T]\), including \(t=0\). That form is false for every \(\nu\):
\begin{lstlisting}
theorem old_periodicDirectionDepletionBound_unsatisfiable (nu : Real) :
    Not (PeriodicDirectionDepletionBoundOriginal nu)
\end{lstlisting}
(\texttt{HighSetDepletionAndBKM.\hspace{0pt}lean}, lines 735--736, first proved in \texttt{afaab80}). The witness is the path that is zero on \([0,\infty)\) and a linear shear with nonzero vorticity for \(t<0\). Every implication that assumed the old form was vacuously true. The corrected one-sided clause is not refuted by values off \([0,T]\) (\texttt{jumpPath\_\hspace{0pt}oneSided\_\hspace{0pt}depletion\_\hspace{0pt}clause}).

The fefferman file also defines a periodic closing bound on the closed interval \([0,T]\) (\texttt{PeriodicTetherClosingBound}, same file, lines 269--276), and proves that the one-sided depletion bound implies that ceiling (\texttt{periodicTetherClosingBound\_\hspace{0pt}of\_\hspace{0pt}directionDepletion}).

\subsection[The same names in NS\_Millennium\_Proof (commit aa9d28c)]{The same names in \texttt{NS\_Millennium\_Proof} (commit \texttt{aa9d28c})}

\textit{Source:} \texttt{BenFrohman/\hspace{0pt}NS\_\hspace{0pt}Millennium\_\hspace{0pt}Proof}, commit \texttt{aa9d28c}. This repository uses some of the same identifiers with \textbf{different meanings}. That fact is recorded, not repaired, here.

{\sloppy\begin{itemize}[leftmargin=1.4em]
\item \texttt{FrohmanianTetherBound} and \texttt{FrohmanianTetherRate} in \texttt{NS\_\hspace{0pt}Millennium\_\hspace{0pt}Proof/\hspace{0pt}Modules/\hspace{0pt}SymplecticTether.\hspace{0pt}lean}, lines 728--734 and 740--751, match the sdiff physical-class definitions above, including the two-sided \texttt{DifferentiableAt} in the rate.
\item \texttt{PeriodicTetherClosingBound} in \texttt{NS\_\hspace{0pt}Millennium\_\hspace{0pt}Proof/\hspace{0pt}Modules/\hspace{0pt}ClayStatements.\hspace{0pt}lean}, line 568, matches the sdiff periodic ceiling (solutions on \([0,T)\)).
\item \texttt{PeriodicDirectionDepletionBound} in \texttt{NS\_\hspace{0pt}Millennium\_\hspace{0pt}Proof/\hspace{0pt}Modules/\hspace{0pt}HighSetDepletionAndBKM.\hspace{0pt}lean}, lines 195--204, is \textbf{not} the cubic rate. It is a uniform envelope
\[
\|\omega(t)\|_{\infty}\le \max\bigl(\|\omega_{0}\|_{\infty},\, C/\kappa_{\mathrm{dir}}\bigr)
\]
along periodic solutions on \([0,T)\). The file itself says the accurate public name is \texttt{UniformVorticityEnvelopeBound} (line~208), and that the definition ``assumes the ceiling itself; it does NOT assert or derive a differential inequality.''
\item A separate predicate \texttt{CubicVorticityDifferentialBound} (lines 212--224) is a compact-interval, two-sided \texttt{DifferentiableAt} rate, declared as ``an additional hypothesis, NOT a theorem about Navier--Stokes solutions.'' It has not been updated to the fefferman one-sided form.
\end{itemize}}

\section{Known refutations and vacuities}
\label{sec:refutations}

These are kernel theorems, not informal caveats.
{\sloppy\begin{itemize}[leftmargin=1.4em]
\item \texttt{TetherRateCounterexample.\hspace{0pt}not\_\hspace{0pt}forall\_\hspace{0pt}frohmanianTetherRate} and \texttt{not\_\hspace{0pt}frohmanianTetherRate\_\hspace{0pt}U} (sdiff \texttt{FrohmanianTether/\hspace{0pt}TetherRateCounterexample.\hspace{0pt}lean}, lines 423--430; the same names in \texttt{NS\_\hspace{0pt}Millennium\_\hspace{0pt}Proof/\hspace{0pt}Modules/\hspace{0pt}TetherRateCounterexample.\hspace{0pt}lean}, lines 623--630). The universal quantification ``for every datum, the physical-class rate holds'' is false. A steady finite-energy field \(U\) with a nowhere-differentiable (Takagi) component ``solves'' \texttt{NS\_PDE} because that class does not demand spatial differentiability and Mathlib's \texttt{fderiv} is zero at non-differentiable points, while \(\|\omega(U)\|\) is unbounded, violating the \texttt{BddAbove} conjunct. The bound is likewise refuted: \texttt{not\_\hspace{0pt}frohmanianTetherBound\_\hspace{0pt}U}.
\item \texttt{not\_\hspace{0pt}forall\_\hspace{0pt}localTetherRateV014} (sdiff \texttt{FrohmanianTether/\hspace{0pt}TetherDefinitionAudit.\hspace{0pt}lean}, line 241, via the manuscript stub \texttt{manuscript/\hspace{0pt}lean/\hspace{0pt}not\_\hspace{0pt}forall\_\hspace{0pt}localTetherRateV014.\hspace{0pt}lean}). The earlier two-sided compact-interval rate is not a theorem for every datum.
\item \texttt{TetherRatePeriodicVacuity}: on nonzero continuous \(\mathbb{Z}^{3}\)-periodic data, \texttt{FrohmanianTetherRate} and \texttt{FrohmanianTetherBound} hold only vacuously, because such data are not square-integrable on the model space used by \texttt{IsPhysicalNSSolution}. This is one precise form of the \(\Tthree\) versus \(\R^{3}\) (or non-\(L^{2}\) periodic) mismatch. The periodic Clay~(B) hypotheses above do not use that physical-class filter.
\item \texttt{old\_\hspace{0pt}periodicDirectionDepletionBound\_\hspace{0pt}unsatisfiable} (fefferman, commit \texttt{afaab80}, retained on main at \texttt{3512caa}): the pre-v1.5.0 two-sided depletion hypothesis is false for every \(\nu\).
\item Zero data are a special case, not a general proof: \texttt{frohmanianTetherRate\_zero} and \texttt{frohmanianTetherBound\_zero} hold unconditionally after the physical-class repair at the zero field.
\end{itemize}}
{\sloppy The repaired statement on Fefferman's smooth class, \texttt{FeffermanClayA.\allowbreak FrohmanianTether\allowbreak RateSmooth}, is recorded in~\S\ref{sec:rate-smooth}. It is open for general data; the refutation of the unrestricted rate is why the statement was repaired. It is not proved here.}

The tether-regularity record \doi{10.5281/zenodo.24144856} separately states that its \texttt{T3} is really \(\R^{3}\), so its regularity theorem says nothing about periodic flows, and that its P5 claim is withdrawn. Those corrections stand on that record; they are not reopened or strengthened here.

\section{Conditional implications}

In every case the conclusion is a proved implication. The named hypotheses remain open, with one exception: periodic BKM continuation in the sdiff library after pull request~\#5, which is a proved theorem for every \(\nu>0\). The \(\R^{3}\) companion \texttt{R3BKMContinuation} is not that theorem; it remains a hypothesis (see~\S\ref{sec:rate-smooth}).

\subsection{sdiff, after the BKM pillar (commit \texttt{4aba751})}

\texttt{FrohmanianTether/\hspace{0pt}Pillars/\hspace{0pt}BKMContinuation.\hspace{0pt}lean}, lines 144 and 170--178:
\begin{lstlisting}
theorem periodicBKMContinuation_of_classical {nu : Real} (hnu : 0 < nu) :
    PeriodicBKMContinuation nu

theorem clay_B_exact_of_periodic_tether_closing_bound
    (hclose : forall nu > (0 : Real), PeriodicTetherClosingBound nu) :
    Clay_B_exact

theorem clay_B_exact_unique_of_periodic_tether_closing_bound
    (hclose : forall nu > (0 : Real), PeriodicTetherClosingBound nu) :
    Clay_B_exact_unique
\end{lstlisting}
The only remaining hypothesis is the open periodic closing bound. Uniqueness of the global periodic velocity is proved unconditionally (\texttt{periodicSmoothSolution\_unique} in \texttt{FrohmanianTether/\hspace{0pt}Pillars/\hspace{0pt}Uniqueness.\hspace{0pt}lean}); \texttt{Clay\_\hspace{0pt}B\_\hspace{0pt}exact\_\hspace{0pt}unique} is therefore equivalent to \texttt{Clay\_\hspace{0pt}B\_\hspace{0pt}exact} (\texttt{clay\_\hspace{0pt}B\_\hspace{0pt}exact\_\hspace{0pt}unique\_\hspace{0pt}iff}). Both sides remain open as statements about nature. The file states, and this note repeats, that the closing bound is, by the Beale--Kato--Majda criterion, as hard as Statement~(B) itself. This is not a proof of Statement~(B).

The older form \texttt{clay\_\hspace{0pt}B\_\hspace{0pt}exact\_\hspace{0pt}of\_\hspace{0pt}periodic\_\hspace{0pt}tether\_\hspace{0pt}bound} (\texttt{ClayStatements.\hspace{0pt}lean}, lines 599--601) still takes both \texttt{hclose} and \texttt{hBKM}; the pillar discharges \texttt{hBKM} by supplying \texttt{periodicBKMContinuation\_\hspace{0pt}of\_\hspace{0pt}classical}.

The library also has a physical-class implication \texttt{frohmanian\_\hspace{0pt}conditional\_\hspace{0pt}regularity\_\hspace{0pt}of\_\hspace{0pt}closing\_\hspace{0pt}bound} (\texttt{ClayStatements.\hspace{0pt}lean}, line 159), which still takes Kato, BKM, uniqueness, and \texttt{LocalTetherClosingBound} as binders, and \texttt{conditional\_\hspace{0pt}regularity\_\hspace{0pt}of\_\hspace{0pt}open\_\hspace{0pt}hyps}. Those are not Statement~(B).

\subsection{sdiff, Fefferman Problem~(A) reduction (same commit)}
\label{sec:clay-a-reduction}

{\sloppy \texttt{FrohmanianTether/\hspace{0pt}FeffermanClayA.\hspace{0pt}lean}, lines 121--125, with the literal conclusion \texttt{Clay\_\hspace{0pt}A\_\hspace{0pt}Fefferman} at lines 55--59:}
\begin{lstlisting}
@[expose] public def Clay_A_Fefferman : Prop :=
  forall (nu : Real) (u0 : VelocityField),
    0 < nu -> ContDiff Real (top : NatInfinity) u0 ->
      (forall x, div u0 x = 0) -> FeffermanDecay u0 ->
        exists (u : TimeDependentVelocity)
          (p : TimeDependentPressure),
          FeffermanGlobalSolution nu u0 u p

public theorem clay_A_fefferman_of_frohmanianTetherRateSmooth
    (hrate : forall nu > (0 : Real), forall u0 : VelocityField,
      ContDiff Real (top : NatInfinity) u0 ->
        (forall x, div u0 x = 0) -> FeffermanDecay u0 ->
          FrohmanianTetherRateSmooth u0 nu)
    (hBKM : forall nu > (0 : Real), R3BKMContinuation nu) :
    Clay_A_Fefferman
\end{lstlisting}
In words, writing ``Fefferman data'' for smooth, divergence-free fields on \(\R^{3}\) that satisfy Fefferman's decay~(4):
\begin{align*}
&\bigl(\forall\nu>0,\;\forall\text{ Fefferman data }u_{0},\;
  \texttt{FrohmanianTether\allowbreak RateSmooth}\,u_{0}\,\nu\bigr) \\
&\qquad\to\;
\bigl(\forall\nu>0,\;\texttt{R3BKMContinuation}\,\nu\bigr)
\;\to\;
\texttt{Clay\_A\_Fefferman}.
\end{align*}
The proof applies the already-formalized comparison \texttt{ceiling\_of\_rate} (lines 105--108), which turns the cubic rate into a uniform ceiling \(\max\bigl(M(u_{0}),C/\kappa\bigr)\) on every local Fefferman-class solution, and then invokes the \(\R^{3}\) continuation hypothesis. Both binders remain open for general data; they are shown only at \(u_{0}=0\). This is not a proof of Fefferman's Problem~(A), and it is not a Clay claim.

Once proofs of those two named hypotheses exist, the already-formalized implication gives \texttt{Clay\_A\_Fefferman} directly. No such proofs are offered here, and no derivation of the repaired rate from the unmodified vorticity equation is supplied.

\subsection{fefferman v1.5.0 (commit \texttt{3512caa})}

\texttt{FeffermanClayNS/\hspace{0pt}MainTheorem.\hspace{0pt}lean}, lines 35--38:
\begin{lstlisting}
theorem clay_reduces_to_direction_curvature_bound :
    forall nu : Real, 0 < nu -> PeriodicBKMContinuation nu ->
      (DirectionCurvatureBound nu -> Clay_B_exact nu)
\end{lstlisting}
Here \texttt{DirectionCurvatureBound} is the corrected one-sided \texttt{PeriodicDirectionDepletionBound}, which is open. \texttt{PeriodicBKMContinuation} remains an explicit hypothesis in this repository (classical BKM, not formalized in this file). The converse is not proved.

\subsection[NS\_Millennium\_Proof v2.5.0 (commit aa9d28c)]{\texttt{NS\_Millennium\_Proof} v2.5.0 (commit \texttt{aa9d28c})}

\texttt{NS\_\hspace{0pt}Millennium\_\hspace{0pt}Proof/\hspace{0pt}Modules/\hspace{0pt}MainTheorem.\hspace{0pt}lean}, lines 30--40:
\begin{lstlisting}
theorem clay_reduces_to_direction_curvature_bound
    (hBKM : forall nu > (0 : Real), PeriodicBKMContinuation nu) :
    (forall nu > (0 : Real), DirectionCurvatureBound nu) -> Clay_B_exact

theorem clay_reduces_to_uniform_vorticity_envelope
    (hBKM : forall nu > (0 : Real), PeriodicBKMContinuation nu) :
    (forall nu > (0 : Real), UniformVorticityEnvelopeBound nu) -> Clay_B_exact
\end{lstlisting}
The two theorems are the same implication: \texttt{DirectionCurvatureBound} is an alias of the envelope \texttt{UniformVorticityEnvelopeBound}, not of a cubic rate. Both the envelope and classical BKM continuation remain hypotheses. No general PDE estimate is supplied. The periodic-bound form \texttt{clay\_\hspace{0pt}B\_\hspace{0pt}exact\_\hspace{0pt}of\_\hspace{0pt}periodic\_\hspace{0pt}tether\_\hspace{0pt}bound} (\texttt{ClayStatements.\hspace{0pt}lean}, lines 592--594) likewise still takes both \texttt{hclose} and \texttt{hBKM}.

\section{Honest scope}

{\sloppy\begin{itemize}[leftmargin=1.4em]
\item The cubic rate and the uniform vorticity closing bound are \textbf{open} for general smooth, divergence-free, finite-energy (or periodic) data. The repaired rate \texttt{FrohmanianTether\allowbreak RateSmooth} is likewise open for general Fefferman data.
\item By Beale--Kato--Majda, a time-uniform \(\|\omega\|_{\infty}\) bound along every local periodic solution, for all such data, is as hard as Statement~(B). The gaps T1--T6 of 22835069 are that same missing a priori bound under different names.
\item The formalizations prove implications, not Statement~(B), not Fefferman's Problem~(A), and not the Clay Millennium problem. In particular, \texttt{clay\_\hspace{0pt}A\_\hspace{0pt}fefferman\_\hspace{0pt}of\_\hspace{0pt}frohmanianTether\allowbreak RateSmooth} would yield \texttt{Clay\_\hspace{0pt}A\_\hspace{0pt}Fefferman} only after proofs of both of its binders exist.
\item \texttt{R3BKMContinuation} (classical continuation on \(\R^{3}\)) remains a hypothesis; it is not the proved periodic theorem \texttt{periodicBKMContinuation\_\hspace{0pt}of\_\hspace{0pt}classical}.
\item Known refutations of the physical-class rate and of the old two-sided depletion hypothesis are theorems; they are not side remarks.
\item Where a repository's \texttt{T3} is really \(\R^{3}\), or where a physical-class filter excludes genuine periodic data, a theorem on that class is silent about Fefferman's periodic Statement~(B).
\item The two-point strain pairing is not free. An \(L^{\infty}\) bound on \(S\) on the whole torus is emptiness of the singular set, not \(\mathbb{P}^{1}\)-nullity.
\end{itemize}}

\section{Reproducibility}

No \texttt{lake build} was run for this note: a lake build was already running on the shared machine, and a new Mathlib cache was not downloaded. Axiom output is reused from existing CI and Zenodo artifacts. All three repositories pin
\begin{center}
toolchain \texttt{leanprover/\hspace{0pt}lean4:v4.\hspace{0pt}30.\hspace{0pt}0-rc1}, \qquad
Mathlib \texttt{0c154d67103f74be3a0f2c509f72ccbf5be9f2a7}.
\end{center}

\paragraph{sdiff (commit \texttt{4aba751}).}
CI: \url{https://github.com/BenFrohman/frohmanian-tether-sdiff/actions/runs/43105941627} (green). The artifact \texttt{axioms-pull\_request-43105941627} contains 1{,}882 \texttt{\#print axioms} records, of which 1{,}868 uniquely parsed names; every recorded axiom set is contained in \(\{\texttt{propext},\texttt{Classical.choice},\texttt{Quot.sound}\}\). This includes \texttt{periodicBKMContinuation\_\hspace{0pt}of\_\hspace{0pt}classical}, \texttt{clay\_\hspace{0pt}B\_\hspace{0pt}exact\_\hspace{0pt}unique\_\hspace{0pt}of\_\hspace{0pt}periodic\_\hspace{0pt}tether\_\hspace{0pt}closing\_\hspace{0pt}bound}, \texttt{not\_\hspace{0pt}forall\_\hspace{0pt}frohmanianTetherRate}, and \texttt{not\_\hspace{0pt}forall\_\hspace{0pt}localTetherRateV014}. One-line rebuild, from a clone at that commit, after \texttt{elan} has the pinned toolchain:
\begin{lstlisting}
lake build && lake env lean scripts/AxiomAudit.lean
\end{lstlisting}
or the repository's \texttt{bash scripts/\hspace{0pt}verify.\hspace{0pt}sh}.

\paragraph{fefferman (commit \texttt{3512caa}, Lean correction in \texttt{afaab80}).}
CI: \url{https://github.com/BenFrohman/fefferman-clay-navier-stokes/actions/runs/43091413968}. The v1.5.0 README records 236 theorem/lemma declarations in source and 243 names audited by \texttt{AxiomAudit.\hspace{0pt}lean}. The file \texttt{AxiomAudit\_\hspace{0pt}output.\hspace{0pt}txt} from \doi{10.5281/zenodo.24144846} contains 243 axiom lines, all \(\{\texttt{propext},\texttt{Classical.choice},\texttt{Quot.sound}\}\), including \texttt{old\_\hspace{0pt}periodicDirectionDepletionBound\_\hspace{0pt}unsatisfiable}. Rebuild:
\begin{lstlisting}
lake build && lake env lean scripts/AxiomAudit.lean
\end{lstlisting}

\paragraph{\texttt{NS\_Millennium\_Proof} (commit \texttt{aa9d28c}).}
Lean CI: \url{https://github.com/BenFrohman/NS\_Millennium\_Proof/actions/runs/43103668767} (green). Pull request~\#29, merged at this commit, changed the paper title and scope wording, not the Lean. The axiom excerpt deposited with \doi{10.5281/zenodo.24144853} (CI run 43047743968) shows the standard three axioms on the audited names, including \texttt{frohmanianTetherRate\_zero} and \texttt{not\_\hspace{0pt}forall\_\hspace{0pt}frohmanianTetherRate}. Rebuild:
\begin{lstlisting}
lake build && lake env lean scripts/AxiomAudit.lean
\end{lstlisting}

\paragraph{This note.}
The PDF is built from \texttt{main.tex} by
\begin{lstlisting}
tectonic main.tex
\end{lstlisting}
The original published cover PDF of 22835069 is included unchanged as \texttt{00\_\hspace{0pt}COVER\_\hspace{0pt}Official\_\hspace{0pt}Proposed\_\hspace{0pt}Theorem.\hspace{0pt}pdf}.

\section*{Copyright and license}

Copyright \(\copyright\) 2026 Benjamin Stanley Frohman. All rights reserved. Released under the Apache License, Version 2.0 (SPDX: Apache-2.0). Author: Frohman, Benjamin Stanley (ORCID 0009-0006-7068-3718; Independent Researcher, Texas, United States). Contact: \texttt{frohmanbenjamin@gmail.com}.

\end{document}
